\section{The Agent Boundary}\label{sec:host}

An agent's tool calls reach the file system, the network and the key keeper
through the agent host (\file{crates/agent-host}), which checks each call
against the agent's grant before acting on it. The host is an in-process
library, not a sandbox for the agent: nothing in it confines the agent's
process, which is why its being the agent's only path to those resources is
an assumption about the deployment\claim{sole-path}, and its \code{call}
takes the caller's agent id as given, so a deployment must bind each agent's
tool-call channel to its id\claim{agent-identity}. Process sandboxes such as
nono~\cite{nono} or a verified runtime such as WaVe~\cite{wave} are what
confine an agent; the host decides what a confined agent's requests may do.

The host has two interfaces. The admin interface issues a root agent with a
grant and revokes an agent together with every agent under
it\claim{host-revocation}. The agent interface takes five tool calls:
\code{read\_file}, \code{write\_file}, \code{fetch}, \code{sign} and
\code{delegate}. None reads or edits a grant\claim{host-no-grant-tool};
\code{delegate} starts a sub-agent whose grant is the caller's caveats
followed by the added ones, so every child allows a subset of its
parent\claim{host-delegation}.

\subsection{The evaluator as a component}

The verdict comes from \file{crates/grant-component}: \code{grant-eval}
compiled for \code{wasm32v1-none}, the target the grant signer uses, and
wrapped as a WebAssembly component behind a WIT interface with one interface
per vocabulary, each exporting \code{allows} and \code{allows-all}. The
component imports nothing and keeps no state between calls, so it can compute
only on the grant and action the host passes in. The host builds the component
at compile time, instantiates it in Wasmtime with an empty linker, and passes
the grant and the action on every call. Since WIT has no 128-bit integer, a
chain \code{i128} crosses the boundary as a pair of 64-bit halves. Tests check
that the component returns the native evaluator's verdict on every grant of up
to three caveats from a small universe per vocabulary, and for chain on every
authorization tree of up to two contexts\claim{component-agrees}; that
Wasmtime and the bindings carry values across faithfully is
trusted\claim{runtime-trusted}.

The evaluator is proven total and pure, so the component does not defend
against the evaluator's logic. It does two other things. It runs the build the
chain runs: the same source for the same \code{wasm32v1-none} target, so the
host's verdicts come from the artifact class the grant signer deploys, not
from a separate native build. And it confines that build, which is the output
of an unverified toolchain, to computing a verdict: a component that imports
nothing can touch no file, socket or key whatever its compiled code does. The
price is the boundary crossing measured in \S\ref{sec:eval}.

\subsection{Building actions}

The wiring that turns a tool call into an action is the trusted part of the
host, kept to three modules.

\paragraph{Files.} A requested path must be absolute. Each existing prefix is
canonicalized with \code{realpath}, which follows symlinks, applies \code{..}
to the directory a symlink led to, and returns the names stored on disk;
components past the last existing one are appended as given. Grant prefixes
are resolved the same way when a grant is issued or delegated. On a
case-insensitive file system the check therefore sees the stored name, and a
differently cased path is allowed exactly when the stored path
is\claim{host-case}. The host checks the canonical path, opens the file with
\code{O\_NOFOLLOW} on its last component, and checks the path of the opened
descriptor again (\code{F\_GETPATH} on macOS, \code{/proc/self/fd} on Linux)
before any data moves through that descriptor, so a component swapped for a
symlink between check and open is caught\claim{host-swap}. A new file is
created under its directory's checked descriptor with
\code{O\_CREAT|O\_EXCL|O\_NOFOLLOW}. Neither \code{..} nor a symlink reaches
outside the grant\claim{host-dotdot}\claim{host-symlink}. Three cases are not
covered. A file renamed or relinked by another process after the descriptor's
path was checked is read or written all the same, since the descriptor follows
the inode, not the path. A hard link inside the grant to a file outside it is
allowed, since the vocabulary names paths, not files. And an error while
resolving is reported before the check, so it can reveal that something
exists at a path the grant does not cover. The agent's own tool calls create
neither renames nor links.

\paragraph{Network.} The action is the host name the agent requested, as the
URL parser normalizes it (lower-case ASCII, IDNA names in \code{xn--} form), with
its port and scheme. After the check, the host's client resolves the name and
connects directly; the agent never supplies an address, and what a name
resolves to is trusted\claim{os-wiring}. A denied host, port or scheme is
refused before any connection\claim{host-network-denied}. The client follows
no redirects itself: each \code{Location} is resolved against the current URL
and checked as a new action before the next request, up to five hops, so a
redirect to a disallowed host is refused without contacting
it\claim{host-redirect}. The vocabulary binds names, not addresses: an
allowed name that resolves, or is rebound, to a loopback, link-local or cloud
metadata address is reached, and the HTTP method and path are not part of the
action. A caveat on the resolved address's class, or refusing private ranges
unless named, would close most of the first.

\paragraph{Signing.} The agent names a key id, a network passphrase and a
payload, and holds no key. A grant carries a chain caveat list beside its
signing list, parsed and mapped by the relay's code. The host checks the
signing action (key, network, payload kind) first; for a transaction or an
authorization entry it then decodes the payload and checks every context the
signature authorizes against the chain grant with \code{allows\_all}, and
only then calls the key keeper, so a denied request never reaches
it\claim{host-signing-before-keeper}. For a transaction the contexts are its
host function and every context of every authorization entry it carries; for
an authorization preimage, the account's own or a delegated signer's, they
are every context of its invocation tree\claim{host-signing-chain}. The host
signs only what the chain vocabulary describes: an envelope with exactly one
\code{InvokeHostFunction} operation, and no fee bump, second operation,
other operation or contract-code upload\claim{host-signing-refuses}. The
keeper is given the encoding of the decoded value the host checked, so the
bytes signed are the bytes checked\claim{host-signs-checked}. Each payload
kind is also bound to its network: a transaction is signed as its hash on the
named network; a preimage is signed only after the host checks that its
network id is the hash of the named passphrase\claim{host-delegated-preimage};
and a message is signed under SEP-53's prefix ``\code{Stellar Signed
Message:\textbackslash n}''~\cite{sep53}, which keeps a message signature from
being a transaction or authorization-entry signature. A message signature is
not bound to any network, so for messages a \code{networkIn} caveat restricts
only the label the agent supplies.

No caveat covers a transaction's fee, sequence number, time bounds, memo or
source account, an authorization entry's nonce or expiration ledger, or which
account a key id stands for, so an agent allowed to sign transactions can have
its key pay any fee for an allowed call. An empty chain list allows every
contract call, as an empty grant does, so a grant that allows
\code{transaction} or \code{authEntry} signatures should carry chain caveats.
